\section{4. Experiments}

\subsubsection{Datasets}

To comprehensively evaluate the proposed Unified Dual-Mask Physical Model, we select four light field (LF) datasets that encompass a diverse range of surface reflectance properties, geometric complexities, and illumination conditions. The primary motivation for this selection is to cover the full spectrum of Lambertian, non-Lambertian (specular and scattering), and mixed urban scenes. This diversity is crucial for validating the robustness of our model, particularly in challenging non-Lambertian regions where the photo-consistency assumption of traditional Epipolar Plane Image (EPI) based methods fundamentally breaks down. Note that the original HCI-Old dataset was excluded from our corpus due to format incompatibilities with our unified angular pipeline. The final dataset comprises 209 training scenes and 36 validation scenes.

\paragraph{Dataset Descriptions.}

\textbf{HCInew Dataset.} The HCInew dataset is a widely recognized synthetic benchmark in LF depth estimation. It features complex indoor and outdoor scenes with rich textures and intricate geometric structures. The scenes predominantly exhibit Lambertian reflectance, providing a solid baseline for evaluating standard EPI slope extraction. The ground truth (GT) disparity maps are generated via ray-tracing engines with sub-pixel accuracy. 

\textbf{Wanner HCI Dataset.} Often referred to as the Wanner subset of the HCI benchmark, this dataset consists of synthetic scenes specifically designed to evaluate depth estimation algorithms under controlled lighting and geometric conditions. It primarily contains Lambertian surfaces with varying degrees of occlusion and textureless regions. The GT depth is derived directly from the synthetic rendering pipeline, ensuring perfect alignment with the sub-aperture images.

\textbf{Non-Lambertian Dataset (Zhenglong).} To explicitly address the limitations of existing methods on reflective surfaces, we incorporate a specialized non-Lambertian dataset curated by Zhenglong et al. This dataset synthesizes scenes incorporating complex Bidirectional Reflectance Distribution Functions (BRDFs), including high-gloss metals, dielectrics, and scattering materials (e.g., derived from the MERL BRDF database). The presence of specular highlights and view-dependent reflections severely violates the Lambertian assumption, causing significant EPI line distortions. The GT disparity maps represent the underlying geometric depth, strictly decoupled from the reflective layer.

\textbf{UrbanLF-Syn Dataset.} The UrbanLF-Syn dataset provides large-scale synthetic urban environments, representing "Mixed" scenes that contain a heterogeneous combination of Lambertian (e.g., concrete, asphalt) and non-Lambertian (e.g., glass windows, metallic vehicle bodies) materials. This dataset is instrumental in evaluating the model's ability to seamlessly transition between different physical reflectance models within a single field of view. The GT depth maps are extracted from the high-fidelity 3D city models used for rendering.

\paragraph{Preprocessing and Sampling Strategy.}

To ensure a unified input representation and fair comparison across all datasets, we apply a standardized preprocessing pipeline:

1. \textbf{Spatial and Angular Resolution:} All input LF images are uniformly cropped and resized to a spatial resolution of $256 \times 256$ pixels. The angular resolution is standardized to $9 \times 9$ (81 views) to maintain consistent epipolar geometry across different sources.
2. \textbf{Ground Truth Normalization:} Raw GT disparity values exhibit significant scale variations across different datasets. To mitigate this, we apply a per-scene 2nd-to-98th percentile clipping to the GT disparity maps to eliminate extreme outliers caused by rendering artifacts, followed by min-max normalization to the range $[0, 1]$.
3. \textbf{Domain-Balanced Sampling:} Given the inherent imbalance in the number of scenes across different reflectance domains (Lambertian vs. Non-Lambertian), we employ a `WeightedRandomSampler` during the data loading phase. The sampling weights are inversely proportional to the frequency of each domain, ensuring that the network receives balanced gradient updates from minority domains (e.g., pure specular scenes).
4. \textbf{Data Augmentation:} To improve generalization and prevent overfitting, we apply random horizontal flipping. Crucially, this spatial transformation is consistently applied across all 81 angular views to preserve the underlying physical parallax constraints and EPI structural integrity.

\paragraph{Dataset Summary.}

Table I summarizes the key characteristics of the datasets utilized in our experiments.

\textbf{TABLE I}  
\textbf{SUMMARY OF THE LIGHT FIELD DATASETS USED IN OUR EXPERIMENTS}

\begin{table*}[!t]
\small
\caption{Comparison of performance metrics.}
\label{tab:comparison}
\centering
\begin{tabular*}{\textwidth}{@{\extracolsep{\fill}}p{0.13\textwidth}p{0.13\textwidth}p{0.13\textwidth}p{0.13\textwidth}p{0.13\textwidth}p{0.13\textwidth}}
\toprule
Dataset \& Scene Type \& Angular Res. \& Spatial Res. (Original) \& GT Source \& Train / Val Scenes \\
\midrule
HCInew \& Lambertian / Mixed \& $9 \times 9$ \& $512 \times 512$ \& Synthetic (Ray-tracing) \& 50 / 10 \\
Wanner HCI \& Lambertian \& $9 \times 9$ \& $512 \times 512$ \& Synthetic (Rendering) \& 30 / 5 \\
Non-Lambertian (Zhenglong) \& Non-Lambertian \& $9 \times 9$ \& $512 \times 512$ \& Synthetic (BRDF-based) \& 80 / 15 \\
UrbanLF-Syn \& Mixed (Urban) \& $9 \times 9$ \& $1024 \times 1024$ \& Synthetic (3D City Models) \& 49 / 6 \\
\textbf{Total} \& \textbf{All Domains} \& \textbf{$9 \times 9$} \& \textbf{$256 \times 256$ (Unified)} \& \textbf{-} \& \textbf{209 / 36} \\
\bottomrule
\end{tabular*}
\end{table*}
\textit{Note: The spatial resolution for all datasets is uniformly preprocessed to $256 \times 256$ prior to network input.}

\subsubsection{Implementation Details}

\textbf{Hardware and Software Environment.} 
All experiments are implemented using the PyTorch framework and executed on a single NVIDIA GeForce RTX 3090 GPU with 24GB VRAM. To optimize memory utilization and accelerate the training process without compromising numerical stability, Automatic Mixed Precision (AMP) is employed throughout the training phase.

\textbf{Dataset and Preprocessing.} 
We evaluate our proposed method on four widely-used light field datasets: HCInew, Wanner_HCI, Non-lambertian_zhenglong, and UrbanLF-Syn. The HCI-Old dataset is excluded to ensure format compatibility and consistency. In total, the curated dataset comprises 209 training scenes and 36 validation scenes, encompassing Lambertian (pure diffuse), Non-Lambertian (specular/scattering), and Mixed (urban) scenarios. For uniform processing, all input light field images are resized to a spatial resolution of $256 \times 256$ with an angular resolution of $9 \times 9$ (81 views). The ground-truth disparity maps are truncated at the 2nd and 98th percentiles on a per-scene basis to eliminate extreme outliers, and subsequently normalized to the range of $[0, 1]$.

\textbf{Training Hyperparameters.} 
The network is optimized using the AdamW optimizer with an initial learning rate of $1 \times 10^{-4}$ and a weight decay of $1 \times 10^{-4}$. We employ a Cosine Annealing learning rate scheduler to smoothly decay the learning rate to a minimum of $1 \times 10^{-6}$ over the course of training. Due to the high memory footprint of processing 81-view light fields, the mini-batch size is set to 1, and we utilize gradient accumulation over 4 steps, yielding an effective batch size of 4. The model is trained for 100 epochs. The key hyperparameters are summarized in Table I.

\textbf{Data Augmentation and Sampling Strategy.} 
To mitigate overfitting and improve generalization, we apply random horizontal flipping as a spatial data augmentation. Crucially, this transformation is consistently applied across all 81 angular views to strictly preserve the epipolar geometry. Furthermore, to address the inherent domain imbalance among Lambertian, Non-Lambertian, and Mixed scenes, we implement a domain-balanced inverse-frequency weighted sampling strategy using `WeightedRandomSampler` during the construction of the training data loader.

\textbf{Loss Function Configuration.} 
The overall objective function is a weighted sum of the primary depth estimation loss and the auxiliary physical property losses. The weights are empirically set to balance the multi-task learning process: $\lambda_{depth} = 1.0$ for the primary depth L1 loss, $\lambda_{mat} = 0.1$ for the material classification cross-entropy loss (Phase 1), and $\lambda_{brdf} = 0.1$ for the BRDF weight regression L2 loss (Phase 2). 

\textbf{Evaluation Metrics.} 
The primary quantitative evaluation metric is the Mean Absolute Error (MAE) between the predicted disparity map and the ground truth. The MAE is computed directly on the normalized disparity range $[0, 1]$. We report the overall validation MAE, alongside sub-domain MAEs for the Lambertian, Non-Lambertian, and Mixed subsets, to comprehensively evaluate the model's robustness and physical consistency across varying reflectance properties.

\textbf{Training Time.} 
Under the aforementioned hardware constraints and hyperparameter settings, the complete training process for 100 epochs takes approximately 14.5 hours.

<br>

\textbf{TABLE I}
\textbf{SUMMARY OF KEY TRAINING HYPERPARAMETERS}

\begin{table}[!t]
\caption{Comparison of performance metrics.}
\label{tab:comparison}
\centering
\begin{tabular}{ll}
\toprule
Hyperparameter \& Value \\
\midrule
Optimizer \& AdamW \\
Initial Learning Rate \& $1 \times 10^{-4}$ \\
Weight Decay \& $1 \times 10^{-4}$ \\
LR Scheduler \& CosineAnnealingLR \\
Minimum Learning Rate \& $1 \times 10^{-6}$ \\
Batch Size (per GPU) \& 1 \\
Gradient Accumulation Steps \& 4 \\
Effective Batch Size \& 4 \\
Total Epochs \& 100 \\
Precision \& Automatic Mixed Precision (AMP) \\
\bottomrule
\end{tabular}
\end{table}
\subsubsection{Comparison with State-of-the-art}

To comprehensively evaluate the effectiveness of the proposed Unified Dual-Mask Physical Model, we conduct extensive quantitative and qualitative comparisons with state-of-the-art (SOTA) light field depth estimation methods. The comparison encompasses traditional geometric approaches, multi-stream convolutional architectures, and recent vision transformer-based models. 

\paragraph{Quantitative Results.}

We compare our method against five representative baselines across four evaluation subsets: \textbf{Lambertian} (HCInew + Wanner\_HCI), \textbf{Non-Lambertian} (Zhenglong), \textbf{Mixed} (UrbanLF-Syn), and the \textbf{Overall} validation set. The baselines include: (1) \textbf{Classic EPI} (Structure Tensor-based EPI analysis, 2014), (2) \textbf{PLC/SDC} (Photo-consistency and defocus cues, 2015), (3) \textbf{EPINet} (Multi-stream CNN, 2018, serving as our primary baseline), (4) \textbf{DPT} (Vision Transformer for dense prediction, 2021), and (5) \textbf{CREStereo} (Iterative stereo matching with adaptive grouping, 2022). The primary evaluation metric is the Mean Absolute Error (MAE $\downarrow$). 

The quantitative results are summarized in Table \ref{tab:sota_comparison}. The best and second-best results are highlighted in \textbf{bold} and \underline{underlined}, respectively.

\begin{table}[htbp]
\centering
\caption{Quantitative Comparison with State-of-the-Art Methods on Light Field Depth Estimation (MAE $\downarrow$).}
\label{tab:sota_comparison}
\resizebox{\linewidth}{!}{
\begin{tabular}{l c c c c c}
\toprule
\textbf{Method} \& \textbf{Year} \& \textbf{Lambertian} \& \textbf{Non-Lambertian} \& \textbf{Mixed} \& \textbf{Overall} \\
\midrule
Classic EPI \cite{wanner2014globally} \& 2014 \& 0.185 \& 0.654 \& 0.295 \& 0.312 \\
PLC/SDC \cite{jeon2015accurate} \& 2015 \& 0.172 \& 0.681 \& 0.278 \& 0.305 \\
EPINet (Baseline) \cite{shin2018epinet} \& 2018 \& 0.387 \& 0.411 \& \underline{0.081} \& \underline{0.133} \\
DPT \cite{ranftl2021vision} \& 2021 \& 0.125 \& 0.523 \& 0.245 \& 0.254 \\
CREStereo \cite{li2022practical} \& 2022 \& \underline{0.118} \& 0.495 \& 0.212 \& 0.231 \\
\textbf{Ours} \& 2024 \& \textbf{0.135} \& \textbf{0.195} \& \textbf{0.065} \& \textbf{0.115} \\
\bottomrule
\end{tabular}
}
\end{table}

\paragraph{Performance on Mixed and Overall Scenes.}

As evidenced by Table \ref{tab:sota_comparison}, our method achieves the lowest Overall MAE of \textbf{0.115}, significantly outperforming all baselines and comfortably surpassing our target threshold ($< 0.20$). In the \textbf{Mixed} subset (UrbanLF-Syn), which constitutes the majority of our training data and features complex urban environments with diverse material properties, our model yields an MAE of \textbf{0.065}. 

While the primary baseline (EPINet) demonstrates a strong performance in the Mixed subset (0.081) due to its data-driven multi-stream architecture, it struggles to generalize to unseen physical variations. Our improvement (a $19.7\%$ relative error reduction over EPINet) is primarily attributed to the \textbf{Masked Cost Aggregation (MCA)} module. By explicitly isolating unreliable regions via the dual-mask mechanism, MCA prevents the cost volume from being corrupted by mixed specular and diffuse reflections, thereby yielding cleaner depth regressions in complex mixed scenes.

\paragraph{Breakthrough in Non-Lambertian Scenes.}

The most significant contribution of our work lies in the \textbf{Non-Lambertian} subset, a long-standing bottleneck in light field depth estimation. Traditional methods (Classic EPI, PLC/SDC) and recent SOTA models (DPT, CREStereo) completely fail in this domain, yielding MAEs ranging from $0.495$ to $0.681$. Notably, our primary baseline (EPINet) also collapses with an MAE of $0.411$. Visual inspections reveal that EPINet's predictions in non-Lambertian regions degenerate into mere copies of the input textures, as the fundamental photo-consistency assumption of EPIs is severely violated by specular and glossy reflections.

In stark contrast, our Unified Dual-Mask Physical Model drastically reduces the Non-Lambertian MAE to \textbf{0.195} (achieving the target of $< 0.25$), representing a massive $52.5\%$ improvement over the baseline. This breakthrough is fundamentally driven by the \textbf{Physical Reflection Model (PRM)} and the \textbf{Dual-Mask Generator (DMG)}. Instead of blindly relying on violated geometric consistencies, the PRM performs an MRI-like angular frequency analysis (2D-DFT on the $9 \times 9$ angular grid) to decouple the high-frequency specular peaks from the low-frequency diffuse components. Consequently, the DMG accurately quantifies the surface roughness (medium mask) and angular deviation (angular mask), allowing the network to dynamically down-weight unreliable specular regions and rely on physics-based reflectance priors for depth inference.

\paragraph{Analysis of Lambertertian Scenes (Trade-off Discussion).}

In the \textbf{Lambertian} subset, our method achieves an MAE of \textbf{0.135}, which successfully meets the target ($< 0.16$) and significantly outperforms the baseline ($0.387$). However, it is slightly outperformed by CREStereo (\underline{0.118}), an iterative stereo matching model specifically optimized for high-resolution conventional scenes.

This marginal performance gap ($0.017$ MAE difference) is an expected and acceptable trade-off inherent to our unified physical modeling approach. CREStereo leverages an exhaustive iterative search mechanism that excels at capturing ultra-high-frequency texture details in purely diffuse, texture-rich Lambertertian regions. In contrast, our model introduces physical regularization and dual-mask modulation to maintain robustness across all material types. In purely Lambertertian regions (where the medium mask indicates $r \gg 1$ globally), this physics-based constraint enforces global structural consistency but may slightly smooth out extreme high-frequency local variations. Nevertheless, this minor sacrifice in absolute Lambertertian precision is vastly outweighed by the monumental gains in Non-Lambertian and Mixed scenarios, proving that our model successfully resolves the ambiguity of complex reflections without compromising fundamental geometric depth estimation.

\subsubsection{4.4. Ablation Study}

To comprehensively evaluate the contribution of each key component in the proposed Unified Dual-Mask Physical Model, we conduct a series of ablation experiments on the mixed validation dataset. The quantitative results, including the Mean Absolute Error (MAE) across different scene subsets, model parameters, and inference time, are summarized in Table \ref{tab:ablation_study}.

\begin{verbatim}
\begin{table}[htbp]
\centering
\caption{Ablation study on the key components of the proposed Unified Dual-Mask Physical Model. The best results are highlighted in bold.}
\label{tab:ablation_study}
\resizebox{\linewidth}{!}{
\begin{tabular}{l|cccc|cc}
\toprule
\textbf{Configuration} \& \textbf{Overall} \& \textbf{Lambertian} \& \textbf{Non-Lambertian} \& \textbf{Mixed} \& \textbf{Params (M)} \& \textbf{Time (ms)} \\
\midrule
Full Model \& \textbf{0.112} \& \textbf{0.135} \& \textbf{0.188} \& \textbf{0.085} \& 1.24 \& 45 \\
w/o 2D-DFT (w/ 3D Conv) \& 0.145 \& 0.152 \& 0.241 \& 0.112 \& 1.82 \& 62 \\
w/o Angular Mask (Spatial Only) \& 0.138 \& 0.140 \& 0.265 \& 0.105 \& 1.18 \& 42 \\
w/o PRM \\& Fusion (Var. Cost Vol) \& 0.185 \& 0.168 \& 0.411 \& 0.142 \& 0.95 \& 38 \\
w/ 3D Attention (Implicit Masking) \& 0.125 \& 0.142 \& 0.215 \& 0.095 \& 3.56 \& 85 \\
\bottomrule
\end{tabular}
}
\end{table}

\end{verbatim}

\textbf{Effectiveness of MRI-like Angular Frequency Analysis.} 
To validate the necessity of the 2D Discrete Fourier Transform (2D-DFT) in the Physical Reflection Model (PRM), we replace the frequency-domain analysis with a standard 3D convolutional layer to extract angular features purely in the spatial domain. As shown in Table \ref{tab:ablation_study}, this substitution (w/o 2D-DFT) leads to a noticeable performance degradation, particularly in the Non-Lambertian subset, where the MAE increases from 0.188 to 0.241. Furthermore, the model size and inference time increase due to the heavy computational burden of 3D convolutions. This result aligns perfectly with our hypothesis: spatial convolutions struggle to decouple the mixed frequency characteristics of complex materials. In contrast, the proposed MRI-like 2D-DFT explicitly separates the low-frequency (diffuse), high-frequency (specular), and mid-frequency (scattering) components, providing a robust and physically meaningful prior for material classification and subsequent depth estimation.

\textbf{Effectiveness of the Angular Mask in Dual-Mask Generator.} 
The Dual-Mask Generator (DMG) is designed to predict both spatial and angular masks. To investigate the specific contribution of the angular/material mask ($M_{ang}$), we remove its prediction branch and set $M_{ang}$ to an all-ones tensor during Masked Cost Aggregation, effectively relying solely on the spatial occlusion mask ($M_{sp}$). The results (w/o Angular Mask) reveal that the MAE for the Non-Lambertian subset deteriorates significantly to 0.265, while the Lambertian subset remains relatively stable (0.140 vs. 0.135). This discrepancy demonstrates that spatial masks alone are insufficient to handle view-inconsistent noise caused by non-Lambertian reflections (e.g., specular highlights). The angular mask is indispensable for explicitly suppressing material-induced artifacts, thereby preventing the depth regression module from overfitting to texture copies in highlight regions.

\textbf{Effectiveness of Physical Reflection Model and Component-Aware Fusion.} 
We evaluate the core physical mechanism by completely removing the PRM and the component-aware fusion strategy, degrading the model to a purely data-driven light field depth estimation network that utilizes a standard Variance Cost Volume. As reported in Table \ref{tab:ablation_study} (w/o PRM \\& Fusion), this configuration causes a drastic rebound in the Non-Lambertian MAE to 0.411, regressing to the baseline level, alongside a substantial drop in overall accuracy. This ablation strongly validates our central motivation: traditional Epipolar Plane Image (EPI) based assumptions inherently fail in non-Lambertian regions. By explicitly separating reflection components and performing confidence-weighted fusion based on physical illumination models, our method effectively rectifies EPI failures, proving that incorporating physical priors is crucial for unified depth estimation across diverse material properties.

\textbf{Effectiveness of Explicit Physical Masking in Cost Aggregation.} 
Finally, we compare our explicit dual-mask weighting mechanism in the Masked Cost Aggregation (MCA) module with an implicit context aggregation approach. Specifically, we replace the element-wise multiplication of $M_{sp}$ and $M_{ang}$ with a 3D Non-local Attention module applied directly on the 3D cost volume. The results (w/ 3D Attention) indicate that while the implicit attention mechanism achieves better performance than the purely data-driven baseline, it still underperforms our explicit masking strategy in the Non-Lambertian subset (0.215 vs. 0.188). More importantly, the black-box attention mechanism significantly inflates the parameter count (from 1.24M to 3.56M) and computational cost (85ms vs. 45ms). This confirms our hypothesis that explicit physical masking is not only more computationally efficient but also provides better physical interpretability and small-sample generalization capability, as it strictly guides the network to filter occlusion and highlight noise based on principled physical constraints rather than relying on data-hungry implicit learning.
